\exnumber{04}
\extitle{UX I - Ex 5 - Problem Statement}
\exfiles{ex04/PDF, named Name\_FirstName\_DeliverableName\_Date\_VersionNumber}
\exauthorize{Those named in the exercise}
\makeheaderfiles

\textbf{Instructions:}

Based on the insights from the analytics and user interviews, comment the scope of possible problem statements.

\textbf{Deliverables:}

\begin{itemize}
\item On a PDF document, write down one problem statement that summarizes the main problem that Pierre has when he runs.
\item Upload the PDF document on Github, labeled as follows: "Name\_FirstName\_DeliverableName\_Date\_VersionNumber". For example, it can be labeled as "Doe\_John\_MidFiPrototype\_05242024\_V1".
\item Add a title within the document.
\item Pay attention to the global aspect of the document. It must be clear, simple and easy to read. You can get inspiration on canva.com but don't overload your design with too much details!
\end{itemize}

\textbf{Tips:}

\begin{itemize}
\item At this stage, it's important to focus on the problem, not on the solution! You can start tackling the problem after (re)defining it.
\item Several problem statements can be possible for each persona and user journey, but for this exercise, we are guiding you towards one.
\item Examples of correct problem statements: 
\begin{itemize}
\item How might we help Pierre (or your persona) measure and prepare his hydration before a race?
\item How might we make sure Pierre doesn't get dehydrated?
\item How might we make sure Pierre drinks the cups that are offered during the race?
\item How might we provide Pierre some drink that contains everything he needs to avoid dehydration?
\item How might make sure Pierre doesn't get drunk the night before the race?
\end{itemize}
\item Examples of wrong problem statements: 
\begin{itemize}
\item How might we make sure Pierre doesn't get dehydrated?
\begin{itemize}
\item Too broad  --- > The problem statement is out of the scope of running, so the solutions cannot be integrated in the specific context of running.
\end{itemize}
\item How might we make sure Pierre drinks the cups that are offered during the race?
\begin{itemize}
\item Too narrow  --- > There are plenty of options to avoid dehydration without drinking the cups that are offered during the race.
\end{itemize}
\item How might we provide Pierre with a drink that contains everything he needs to avoid dehydration?
\begin{itemize}
\item Focused on the solution, not the problem  --- > The problem statement is supposed to rephrase the brief and take into account the insights from the previous phase.
\end{itemize}
\item How might make sure Pierre doesn't get drunk the night before the race?
\begin{itemize}
\item Out of scope  --- > This is not about dehydration, nor about running.
\end{itemize}
\end{itemize}
\end{itemize}

\textbf{Resources:}

\begin{itemize}
\item Stuck on a problem? Just ask how might we \url{https://relab.academy/design-thinking/stuck-on-a-problem-just-ask-how-might-we/}
\item Problem statement \url{https://www.indeed.com/career-advice/career-development/what-is-a-problem-statement}
\item Using 'How Might We' Questions to Ideative on the Right Problems \url{https://www.nngroup.com/articles/how-might-we-questions/}
\end{itemize}

\textbf{Quote:}

"If I had an hour to solve a problem, I'd spend 55 minutes thinking about the problem and 5 minutes thinking about solutions." Albert Einstein

